\subsection{Invitation}
I would like to invite you to participate in this research project which forms part of my Master of Sciences. Before you decide whether you want to take part, it is important for you to understand why the research is being done and what your participation will involve. Please take time to read the following information carefully and discuss it with others if you wish. Ask me if there is anything that is not clear or if you would like more information.
\subsection{What is the purpose of the project?}
This project aims to see if map data, also called geolocation data, can help show how people take part in daily life. Map data can show where someone is compared to nearby places (A relative location e.g. 5km away rather than global co-ordinates), or what kind of place it is (like a park, shop, residential address or restaurant), how often that kind of place is visited in say one week, how long a particular visit lasts, how much travelling is done to get there, and what kind of transport was used (car, public transport, walking etc.). The idea is to use this information to measure “participation,” which means how much people are involved in different parts of life. 

Today, participation is often measured using question sheets. The question sheets take a while to complete and ask people to remember the kinds of things they did, how often they do them normally, and how important doing the thing is to them. Map data might give a new way to measure participation that does not rely on memory or self-scoring how important an activity is to a person. This could help us understand participation better for many different people because it could provided a day-by-day record of how a person has participated without retrospective scoring of that days events being required.

\subsection{``What would be required of me?"}
In broad terms, the requirements for participating in this study are:
\begin{itemize}
    \item Own your own smartphone (IOS 26.0+, Android 8.0+)
    \item Be a healthy adult aged 18-60 with capacity
    \item Have your Smartphone on your person, or to hand for the duration.
    \item Install the research app and leave it running in the background.
    \item Be in the U.K for the duration of data collection
    \item Not void any of the exclusion criteria %link to full form
    \item{Attend an onboarding session online/in-person on \textbf{Day 1}}
    \item{Attend a semi-structures interview online/in-person on \textbf{Day 11+}}
    \item{\textit{Optionally, provide feedback on your experience of being part of the study}}
    \begin{itemize}
        \item{\textbf{Day 1 - 10:} Have your phone about you as you go about your day to day life}
        \item{\textbf{Day 2 - 11:} Populate an app diary confirming/rejecting suggestions about the types of places you visited, the journeys taken and they type of activity you did on the previous day}
        \item{\textbf{Day 11+:} Provide responses to 2 questionnaires during the semi-strucutred interview}
    \end{itemize}

\end{itemize}

\centering{
\textbf{{All of these activities will happen alongside your normal daily routine. You do not need to go to specific places or follow a special schedule to take part in this study, other than the day-by-day events described above. You exact dates fot participating can be arranged to be convenient for you.}}

\textbf{\color{red}{If you are unable or unwilling to meet the criteria above, thanks for reading this far and for your consideration. You don't have to participate and won't be disadvantaged by not participating.}}

\textbf{\color{blue}{If you satisfy the above criteria and would like to participate (or just  be nosey); Please keep readiing to understand more about this study and how it will be carried out.}}

}

\subsection{``What data will I provide and how will it be used"?}
    \begin{itemize}
        \item{\textbf{Personal details:} You will provide your name, age, gender, contact details and ethnicity. This information will be stored in a seperate sheet from all other data. All record of your name and contact details will be destroyed after your exit from the study.}
        \item{\textbf{Maps data:} See below - This data will be used to guess what events you particpated in each day and what type of place you were at.}
        \begin{itemize}
            \item{\textbf{Geolocation data:}} Your time-stamped location relative to starting location (e.g. 400m north-west from start)
            \item{\textbf{Geocoding data:} (The type(s) of place nearest to you visited locations e.g park at 13:12, the name or exact location of the park will never be stored or visible.)}
        \end{itemize}
        \item{\textbf{Diary Responses:} Your responses accepting/rejecting the types of places you visited each day and what activities you did there. This data will act as the truth-record for the activities you particpated in and they type of place you were at.}
        \item{\textbf{Questionnaire Responses:}Your responses to two questionnaire-type measures of participation. These data will be used with your maps and diary data to unerstand which aspects of particpation, as measured by questionnaire, can be correlated with maps data summary measures.}
    \end{itemize}
    {\centering{\textbf{You will remain anonymous in any research outputs.}}}

\subsection{``What if I change my mind about taking part?"'}
You are free to withdraw at any point without having to give a reason. Withdrawing from the project will not affect you in any way. You are able to withdraw your data from the project up until 01/05/2026, after which withdrawal of your data will no longer be possible because the data will have been anonymised or committed to the final report. If you withdraw, the data you have provided will be destroyed.


 